%% To submit your paper:
\documentclass[draft]{agujournal2019}
\usepackage{url} %this package should fix any errors with URLs in refs.
\usepackage{lineno}
\usepackage[inline]{trackchanges} %for better track changes. finalnew option will compile document with changes incorporated.
\usepackage{soul}
\usepackage{amssymb} %for checkmarks in SI table
\usepackage{longtable} %standard multi-page table support
\linenumbers
%%%%%%%
% As of 2018 we recommend use of the TrackChanges package to mark revisions.
% The trackchanges package adds five new LaTeX commands:
%
%  \note[editor]{The note}
%  \annote[editor]{Text to annotate}{The note}
%  \add[editor]{Text to add}
%  \remove[editor]{Text to remove}
%  \change[editor]{Text to remove}{Text to add}
%
% complete documentation is here: http://trackchanges.sourceforge.net/
%%%%%%%

\draftfalse

%% Enter journal name below.
%% Choose from this list of Journals:
%
% JGR: Atmospheres
% JGR: Biogeosciences
% JGR: Earth Surface
% JGR: Oceans
% JGR: Planets
% JGR: Solid Earth
% JGR: Space Physics
% Global Biogeochemical Cycles
% Geophysical Research Letters
% Paleoceanography and Paleoclimatology
% Radio Science
% Reviews of Geophysics
% Tectonics
% Space Weather
% Water Resources Research
% Geochemistry, Geophysics, Geosystems
% Journal of Advances in Modeling Earth Systems (JAMES)
% Earth's Future
% Earth and Space Science
% Geohealth
%
% ie, \journalname{Water Resources Research}

\journalname{JGR: Space Physics}


\begin{document}


\title{GIC Data for Space Weather Event Response Validation \& Evaluation (SWERVE)}


\authors{L. A. Wilkerson\affil{1}, R. S. Weigel\affil{1}, D. Bor\affil{1}, D. Thomas\affil{1}, C. T. Gaunt\affil{2}, P. J. Cilliers, E. J. Oughton\affil{1}}


\affiliation{1}{George Mason University}
\affiliation{2}{University of Cape Town}
%(repeat as many times as is necessary)



\correspondingauthor{L. A. Wilkerson}{lwilker@gmu.edu}



\begin{keypoints}
\item enter point 1 here
\item enter point 2 here
\item enter point 3 here
\end{keypoints}


\begin{abstract}
Geomagnetically induced currents (GICs) are a byproduct of geomagnetic disturbances and have been known to cause disruptions to electric power systems at all latitudes. Access to GIC observations is essential for understanding and assessing GIC risk, including understanding impacts and validating predictive models. However, open GIC data are not centrally accessible, are in a variety of file formats, are often considered proprietary, and can have data quality issues. To address these barriers to GIC research, we have a) assembled a large and unique set of global measured GIC data from both open-access publications and open GIC datasets, b) developed readers for many of the file formats included in the assembled data for ease of access that map the data to a common data structure, and c) built a set of user-selectable algorithms for correcting measured GIC data and rejecting GIC datasets with common issues such as data gaps, low signal, and excessive noise. Several test timeseries were run through these algorithms to ensure their validity, and all thresholds used in the algorithms may be adjusted by the user. 
% make point about data quality being poor
\end{abstract}

\section*{Plain Language Summary}
Geomagnetically induced currents (GICs) present a risk to electric power systems during geomagnetic storms. In order to better understand and more accurately predict the impacts of GICs, observed GIC data must be made accessible to researchers. These data are kept in various places and in various file formats, and they also can have quality issues. This paper addresses these problems by first, collecting open GIC data from around the world; next, putting these data in a common format; and finally, providing an algorithm that corrects common issues with measured GIC data and selects sites based on limits set externally.



\section{Introduction}

Geomagnetic disturbances (GMDs) cause many natural space weather hazards, including geomagnetically induced currents (GICs), a result of electric currents in Earth’s ionosphere and magnetosphere during GMD events (\citeA{Winter2019}; \citeA{Gopalswamy2022}). Solar activity, such as coronal mass ejections, often drives enhancements in ionospheric and magnetospheric currents. These enhancements result in geoelectric field variations on Earth's surface, which in turn cause GICs. GICs, like the geoelectric field, vary in amplitude depending on factors such as geomagnetic latitude and ground conductivity \cite{Wawrzaszek2024}. While stronger GICs are usually found at high latitudes, significant GICs have been observed at middle and low latitudes (\citeA{Liu2018}; \citeA{Caraballo2023}). As such, GICs and their impacts are of interest at all latitudes.

GICs are recognized as a ``potential catastrophic natural hazard" \cite{Lugaz2020} due to their impact on essential infrastructure consisting of long systems of electrical conductors, such as pipelines (\citeA{Khanal2019}; \citeA{Ingham2025}), submarine cables (\citeA{Meloni1983}; \citeA{Chakraborty2022}; \citeA{Boteler2024}), railway signaling and switching systems (\citeA{Love2019}; \citeA{Boteler2020}; \citeA{Patterson2024}), and electrical power transmission networks (\citeA{Gaunt2015}; \citeA{MacManus2025}). 
The impact of GICs on power grids has been of particular interest since the 1989 Hydro-Québec blackout \cite{Bolduc2002}, in addition to GIC-related transformer failures in New Zealand in November 2001 \cite{Marshall2012} and South Africa in October 2003 \cite{Gaunt2007}, both of which notably occurred at middle latitudes. Most recently, a regional power outage in Australia associated with the 19 January 2026 geomagnetic storm left about 8,000 customers without power for a period of around 30~min \cite{Mandow2026}.
These power grid failure events have motivated GIC research globally, with studies in
Canada \cite{Cordell2024}, the US \cite{Wilkerson2026}, Mexico \cite{Caraballo2023}, Brazil \cite{Espinosa2023}, Uruguay \cite{Caraballo2013},  Ethiopia \cite{Liu2018}, South Africa \cite{Ngwira2011}, Ireland \cite{Blake2018}, the UK \cite{Hubert2024}, Spain \cite{Torta2023}, Portugal \cite{AlvesRibeiro2023}, Italy \cite{Pignatiello2026}, Austria \cite{Bailey2022}, Sweden \cite{Rosenqvist2025}, Finland \cite{Waghule2024}, India \cite{Vishnupriya2025}, New Zealand \cite{MacManus2025}, Australia \cite{Waters2025}, Japan \cite{Watari2025}, China \cite{Zhang2020}, and elsewhere \cite{Kelbert2019}. Furthermore, recent socioeconomic models show that power grid disruptions, particularly transformer failures, due to GICs can negatively impact a nation's GDP irregardless of latitude (\citeA{Oughton2017}; \citeA{Oughton2018}; \citeA{Oughton2025}). 

While GIC observations are essential for better understanding, predicting, and mitigating the impacts of GICs on power grids \cite{MacManus2025}, these data are not always openly available due to grid security concerns and most diagnostic data being owned by network operators \cite{Hapgood2016}. Furthermore, open GIC data have not always been cleaned or processed to ensure they are of a usable quality.
For example,  \citeA{Wilkerson2026} found that out of 396 GIC monitors with data openly available through the North American Electric Reliability Corporation (NERC), 355 had poor data quality for the May 2024 ``Gannon" storm. As a result, ensuring the accessibility and usability of measured GIC data that is openly available is of interest to the research community.

In this paper, we present the Space Weather Event Response Validation \& Evaluation (SWERVE) GMD dataset, which includes GIC monitor and supplemental ground magnetometer data during 16 major GMD events since 2000. These data are collected from public data releases and open-access publications globally. Additionally, we present several algorithms for GIC data correction and site selection that correct for baseline offset and large, nonphysical spikes, identify common issues with GIC data (e.g., data gaps, low signal), and flag GIC monitor sites with these issues for exclusion from analysis. 

\section{Data}

\subsection{Data Collection}
There are two primary methods of observing GICs in a given power system: direct observation via Hall effect sensor and indirect observation via differential magnetometer measurement (DMM) \cite{Parry2024}. Hall effect sensors take advantage of the low-frequency nature of GICs by measuring the DC component of the current flowing through the grounded neutral wire at a substation.
However, these sensors are expensive and difficult to maintain, and, as a result, their use globally is limited \cite{Hubert2020}.
DMM uses two magnetometers, one at a remote site perpendicular to the power line that serves as a reference measurement for the variation of the background magnetic field and one placed directly under the power line to measure any excess current flow in addition to the variation of the background magnetic field. The difference between the two measurements yields the magnetic field associated with the GIC, from which the GIC through the line can be inferred using Ampère's law. While DMM is often not used by power grid operators due to the complications of deriving indirect measurements \cite{Marsal2021}, it is often employed by researchers as DMM is cheaper and more flexible as it does not require direct access to substations \cite{Parry2024}.

\subsection{Data Sources}
%NERC, idk if we can include TVA? Talk about mandatory NERC reporting and add details of NERC ERO portal.

NERC is an international regulatory body for the United States, Canada, and Mexico whose stated purpose is ``to assure the effective and efficient reduction of risks to the reliability and security of the grid" \cite{NERC2026_mission}.
NERC, under its transmission planning standards, has published requirements for GMD response, which are dynamic and evolve along with understanding of GMD impacts on electrical power networks (\citeA{NERC2020}; \citeA{NERC2024}).
The Federal Energy Regulatory Commission (FERC) Order No. 830 mandates NERC to collect GIC monitor and magnetometer data and to make such data publicly available in order to support research and analysis of GMD risk \cite{FERC2016}. As such, the NERC Electric Reliability Organization (ERO; \citeA{NERC2026_ERO}) data portal contains reported GIC monitor and ground magnetometer data from various power grid operators across North America for all GMDs with Kp $\geq7$ since May 2013 \cite{NERC2026_reporting}.

Other GIC data was compiled from open-access publications. These data span GIC measurement sites across the globe for a variety of GMD events and provide varying degrees of information during the given GMD. [INCLUDE TABLE W ALL CITATIONS IN APPENDIX]


\subsection{Data Content}
The data collected from NERC includes GIC monitor and ground magnetometer data at NERC sites for 14 storms over North America. The selected storms include all storms since 2020 through June 2025 listed on NERC ERO portal with reported Kp $\geq8$. These storms were selected both for their intensity and recency (earlier listed storms had considerably less sites available). %am currently adding more NERC storms, so this statement is not exactly true
GMD data from both GIC monitors and ground magnetometers is reported to NERC annually by June 30 \cite{NERC2026_reporting}; as such, released datasets from GMDs within the past year may not be complete and are not included in this paper.
Each GIC site has a known geographic latitude, geographic longitude, installation type, and connection. For consistency, sites with connections listed as ``neutral of single-phase transformer" had all GIC values multiplied by a factor of three to match the scaling of the rest of the monitors, which were either listed as ``Common neutral of a 3-phase transformer" or ``Common neutral of three (3) single-phase transformers".
Similarly, each magnetometer site has a known geographic latitude, geographic longitude, and orientation (either geographic or geomagnetic).

The data collected from open-access publications varies in structure and information about the monitor system included in the metadata.

%Include table of storm, \#GIC, \#Mag, maybe some solar wind data to indicate intensity?

\subsection{Data Samples}
Show sample GIC data, stack plots, etc. and discuss

%\subsection{Metadata}
%Describe metadata tables

\section{GIC Timeseries Correction and Selection}


% section on cleaning?? offset, spike removal etc.

\subsection{Correction Filters}

To ensure GIC data quality, we have developed an algorithm to determine potential problems with the GIC data. This algorithm can be divided into two categories: data correction, or correction filters, and site selection, or selection filters.
For each GMD event, a period of heightened GIC activity, called the ``storm time", is defined from visual inspection of the data. 
Our selection algorithm then inspects data from each GIC monitor site available during the storm selectively within given ``storm time" range because some monitors only collect high quality data during periods of high activity. 

First, two correction filters are applied to clean the raw GIC data. The first correction filter detects large, nonphysical spikes and removes them and surrounding points. The second correction filter corrects for GIC timeseries with a baseline offset. If the magnitude of the mean GIC value over the entire dataset is greater than an externally set offset threshold, the mean value is subtracted from the full timeseries.

There are two spike filter options that can be chosen by the user.
The first method of spike detection and removal is a difference filter, which defines a spike as having a sudden, sharp peak that is large compared to the local standard deviation. For each GIC timeseries, the difference between each point was used as a proxy for the first derivative, and the difference between each difference, shifted by one, was used as a proxy for the second derivative. The windows used to compute the local difference, second difference, and standard deviation can be set externally, as can the thresholds used for spike detection. If at a given point the difference and second difference are above their set thresholds while the standard deviation is below its threshold, then the point is marked as a spike.
The second spike detection and removal method uses the median value of a moving window. For each GIC timeseries, a window of given size is moved throughout the duration of the data, and the median value in the window is calculated. The absolute difference between the center point of the window and the median value of the window is then calculated. If this difference is above an externally set threshold, then the point is marked as a spike.

The threshold for what is considered a baseline offset, spike thresholds, and window sizes to remove spikes can all be set externally by the user to make the correction filters as selective as needed. Additionally, there is an option to resample the measured data to a certain sampling rate (default 1~min) to aid in measured-modeled GIC comparison as models typically have a lower data frequency due to the computational cost of getting timeseries results on such a fine scale. This resampling is done after the selection algorithm, described in the next section, has been applied.

\subsection{Selection Criteria}
Once the correction filters have been applied, the selection algorithm checks to see if all  the GIC values positive or negative at the given site, as GIC values should fluctuate around a baseline of 0~A. All other selection criteria presented in this section can have their thresholds changed externally, and each option can be turned off if desired. The selection filters account for five common issues found in the measured GIC data: consistently low values throughout the storm duration, excess levels of noise, large gaps in the data during the storm, low data sampling rate, and constant values during the storm. 

Removing sites with low signal is sometimes a step performed in GIC analysis as very low-amplitude GICs are unlikely to cause thermal impacts on transformer performance; however, the exact threshold at which a transformer is thermally impacted depends not only on the amplitude of the GIC, but also on its duration and the physical characteristics of the transformer such as design and current condition \cite{NERC2017}. Additionally, transformers 
To check for consistent low signal, the selection filter flags all sites that have all values within a set threshold (default $\pm4$~A). 

Noise in the data is defined by comparing standard deviation of the data pre-storm to the standard deviation during storm. If the standard deviation of the data before the storm is greater than a set ratio of the standard deviation during the storm (default $1/4$), then the site is flagged as noisy. The default value is based on visual inspection of GIC data for the May 2024 Gannon storm and should be lowered for weaker storms.

NERC's data reporting requirements advise that device sampling rates should be between 1--10~s, with 1~min being acceptable in cases of equipment limitations \cite{NERC2026_reporting}. However, sampling rates of 1~min and even up to 10~s can lead to the omission of GIC peak values within the 40--70~A range \cite{Trichtchenko2021}. As such, sites with gaps lasting a set time (default 5~min) in data during storm time, including gaps due to spike removal, are also flagged for removal, as well as sites with a sampling rate greater than a set sampling rate (default 1~min) during storm time and sites with constant values during storm time for a set period (default 5~min).

\subsection{Testing and Validation}

To test the correction filters, three test timeseries were developed. The first was a random walk centered on zero with inserted spikes every 3~hr, alternating between positive and negative and increasing in magnitude. The second was a sine wave with an offset baseline (both positive and negative offsets were tested). The third was a combination, a random walk with spikes inserted as before and with an offset. Figure~\ref{} shows examples of these test timeseries before and after the correction filters were applied with default options.

To test the selection algorithm, several test timeseries were created, each related to a specific selection criterion. These test data include: a timeseries with all values within $\pm4$~A, a random timeseries with larger amplitudes during storm time, a timeseries with gaps of various sizes added every 3~hours, a timeseries with a sampling rate of 10~min, and a timeseries with periods of constant values varying in length every 3 hours. Combinations of these tests and the tests for the correction filters were also made to ensure corrected data still passed through the selection filters properly. Examples of these test timeseries are shown in Figure~\ref{}.

Outputs of diagnostic plots, visual/manual validation

\section{Summary}
The SWERVE GMD dataset contains both raw and filtered GIC data from global sources, openly accessible. 
Additionally, code for processing, correcting, and selecting GIC site data is included to address common issues with GIC datasets. While no single filter or rejection criteria will perfectly capture all errors found in GIC data, our approach allows various options to allow testing that conclusions are robust against the filter criteria and their options. We also note that manual inspection is often needed and provide visualization tools to facilitate this step.

%%%%%%%%%%%%%%%%%%%%%%%%%%%%%%%%%%%%%%%%%%%%%%%
%
% DATA SECTION and ACKNOWLEDGMENTS
%
%%%%%%%%%%%%%%%%%%%%%%%%%%%%%%%%%%%%%%%%%%%%%%%

\section*{Open Research Section}
This section MUST contain a statement that describes where the data supporting the conclusions can be obtained. Data cannot be listed as ''Available from authors'' or stored solely in supporting information. Citations to archived data should be included in your reference list. Wiley will publish it as a separate section on the paper’s page. Examples and complete information are here:
https://www.agu.org/Publish with AGU/Publish/Author Resources/Data for Authors


\section*{Conflict of Interest declaration}
The authors declare there are no conflicts of interest for this manuscript.


\acknowledgments
Enter acknowledgments here. This section is to acknowledge funding, thank colleagues, enter any secondary affiliations, and so on.

\section*{Supplemental Information}
%\begin{table}
\begin{longtable}{p{0.106\linewidth}p{0.106\linewidth}p{0.106\linewidth}p{0.106\linewidth}p{0.106\linewidth}p{0.106\linewidth}p{0.106\linewidth}p{0.106\linewidth}p{0.106\linewidth}}
\hline
Author & Location & Events & DMM GIC & Hall GIC & Other Data & Number of GIC sites & GIC data access & Reader in SWERVE? \\ \hline
\endhead
TVA & USA & 2024-05-10, \newline 2023-03-25 &  & $\checkmark$ & measured B, \newline calculated GIC & multiple sites & via request & Yes \\ \hline
NERC & USA & 2013-present &  & $\checkmark$ & measured B & multiple sites & via ERO portal & Yes \\ \hline
SANSA & RSA & multi-event & $\checkmark$ &  &  & multiple sites & via request (Pierre) & Pending \\ \hline
Feng et al 2026 & NZL & 2013-10-02, \newline 2015-03-17, \newline 2015-06-23, \newline 2017-09-08 &  & $\checkmark$ & measured reactive power & 2 sites & NA & No \\ \hline
Wilkerson et al 2026 & USA & 2024-05-10 &  & $\checkmark$ & measured B, \newline calculated GIC, \newline calculated B & 49 sites & https://zenodo.org/records/20090223 & Yes \\ \hline
Belakhovsky et al 2026 & RUS & 2011 - 2022 &  & $\checkmark$ & measured B & 5 sites & http://eurisgic.ru/ (plots only, \newline numerical data exclusive to EURISGIC members) & No \\ \hline
Wang et al 2026 & CHN & 2024-05-10 &  & $\checkmark$ & measured B, \newline calculated GIC & 2 sites & NA & No \\ \hline
Waters et al 2025 & AUS & 2024-05-10, \newline 2024-10-10 &  & $\checkmark$ & measured B, \newline calculated GIC index & 10 sites & NA & No \\ \hline
Parry et al 2025 & CAN & 2023-04-24 &  & $\checkmark$ & measured B, \newline calculated GIC & 2 sites & https://doi.org/10.5281/zenodo.16748607 & Yes \\ \hline
Zhang et al 2025 & CHN & 2012-07-23, \newline 2024-10-10 &  &  & measured B, \newline calculated B, \newline calculated GIC & 25 sites & NA & No \\ \hline
Smith et al 2025 & NZL/UK & 2011 - 2016 &  & $\checkmark$ & measured B, \newline calculated GIC & 24 sites & NA & No \\ \hline
Rosenqvist et al 2025 & SWE & 1999 - 2024 &  & $\checkmark$ & measured B, \newline calculated GIC & 1 site & NA & No \\ \hline
Mac Manus et al 2025 & NZL & 2024-05-10 &  & $\checkmark$ & calculated GIC & 28 sites & NA & No \\ \hline
Amaechi et al 2025 & global & 2023-03-23, \newline 2023-04-23 &  &  & measured B, \newline GIC indices & 15 sites & NA & No \\ \hline
Clilverd et al 2025 & NZL & 2024-05-10 &  & $\checkmark$ & measured B, \newline VLF measurements & 1 site & NA & No \\ \hline
Marsal et al 2025 & ESP & 2021 - 2024 & $\checkmark$ &  & measured B, \newline calculated GIC & 7 sites & https://doi.org/10.34810/data1741 & Pending \\ \hline
Liu et al 2024 & global & 2024-04-23 &  & $\checkmark$ & measured B & 5 sites & NA (see NERC, \newline Cordell et al 2024, \newline FMI, \newline NZL) & Pending \\ \hline
Pratscher et al 2024 & NZL & 2015-03-17, \newline 2017-09-07 &  & $\checkmark$ & measured B, \newline calculated GIC & 4 sites & NA & No \\ \hline
Waghule et al 2024 & FIN & 2013-03-17 &  &  & measured GIC (pipeline), \newline measured B & 1 site & https://space.fmi.fi/gic/man\_ascii/ (FMI) & Pending \\ \hline
Parry et al 2024 & CAN & 2021-10-12 & $\checkmark$ & $\checkmark$ & measured B & 1 site & https://doi.org/10.5281/zenodo.11188122 & Pending \\ \hline
Hubert et al 2024 & GBR & 2018 - 2021 & $\checkmark$ &  & calculated GIC & 12 sites & https://www.bgs.ac.uk/services/ngdc/ & Pending \\ \hline
Smith et al 2024 & NZL & 2001 - 2020 &  & $\checkmark$ & measured B & 22 sites & NA & No \\ \hline
Andreyev et al 2023 & KAZ & 2011-09-26, \newline 2015-06-22, \newline 2016-10-24, \newline 2021-05-12 &  &  & calculated GIC, \newline measured B & 18 sites & NA & No \\ \hline
Ngwira et al 2023 & USA/CAN & 2010 - 2021 &  & $\checkmark$ & measured B & 17 sites & NERC/EPRI & Pending \\ \hline
Caraballo et al 2023 & MEX & Aug-Nov 2021 &  & $\checkmark$ & measured B, \newline calculated GIC & 1 site & https://doi.org/10.17605/OSF.IO/AV9HT “GIC\_assess\_v2:SW Journal Article 2022/datasets” & Pending \\ \hline
Torta et al 2023 & ESP & 1997 - 2022 &  &  & measured B, \newline calculated GIC & 76 sites & NA & No \\ \hline
Alves Ribeiro et al 2023 & PRT & 2010 - 2017 &  & $\checkmark$ & measured B, \newline calculated GIC & 1 site & https://doi.org/10.5281/zenodo.7443212 (GIC statistics), \newline https://doi.org/10.5281/zenodo.7446370 (GIC measurements/estimates) & Pending \\ \hline
Espinosa et al 2023 & BRA & 2015-06-21, \newline 2015-12-19 &  &  & measured B, \newline calculated GIC & 23 sites & NA & No \\ \hline
Ingham et al 2022 & NZL & 2017 - 2020 &  &  & measured GIC (pipeline), \newline measured B & 10 sites & NA & No \\ \hline
Mac Manus et al 2022 & NZL & 1989-03-13, \newline 2003-10-29, \newline 2017-09-07 &  &  & calculated GIC & 33 sites & NA & No \\ \hline
Marshall et al 2022 & AUS & 1998 - 2008 &  &  & GIC indices, \newline measured B & 7 sites & https://doi.org/10.5281/zenodo.7091161 & Pending \\ \hline
Nahayo et al 2022 & RSA & 2003-10-29, \newline 2015-03-17 &  & $\checkmark$ & LDi and LCi indices, \newline measured B & 2 sites & https://zenodo.org/record/7019253 & Pending \\ \hline
Smith et al 2022 & NZL & 2001 - 2016 &  & $\checkmark$ & measured B & 1 site & NA & No \\ \hline
Zhang et al 2022 & CHN & 2012-07-23 &  &  & calculated B, \newline calculated GIC & 54 sites & https://doi.org/10.5281/zenodo.6099318 (code for calculation) & Pending \\ \hline
Bailey et al 2022 & AUT & 1996 - 2022 &  & $\checkmark$ & measured B, \newline calculated GIC & 2 sites & https://doi.org/10.6084/m9.figshare.19102772.v1 (subset of data used) & Pending \\ \hline
Hajra 2022 & FIN & 1999 - 2019 &  &  & measured GIC (pipeline), \newline measured B & 1 site & https://space.fmi.fi/gic/index.php (FMI) & Pending \\ \hline
Rosenqvist et al 2022 & SWE & 2018-05-24, \newline 2000 - 2017 &  & $\checkmark$ & measured B, \newline calculated GIC, \newline calculated B & 1 site & NA & No \\ \hline
Mac Manus et al 2022 & NZL & 2001 - 2020 &  & $\checkmark$ & measured B, \newline calculated GIC & 73 sites & NA & No \\ \hline
Kellerman et al 2021 & USA & 2018 - 2019 &  & $\checkmark$ & measured B, \newline calculated GIC & 7 sites & NERC/EPRI, \newline https://zenodo.org/record/4444068 & Pending \\ \hline
Hughes et al 2021 & USA & 2018 &  & $\checkmark$ & measured B & 9 sites & NERC/EPRI, \newline https://zenodo.org/record/4444068\#.YUoLl9NKj0o & Pending \\ \hline
Albert et al 2021 & AUT & 2017-09-08, \newline 2021-05-12 &  & $\checkmark$ & measured B, \newline calculated GIC & 7 sites & NA & No \\ \hline
Torta et al 2021 & ESP & 2020-08-28 &  & $\checkmark$ & measured B, \newline calculated GIC & 2 sites & NA & No \\ \hline
Trichtchenko 2021 & USA & 2004-07-26, \newline 2001-03-31 &  & $\checkmark$ & measured B & 1 site & EPRI/NERC & Pending \\ \hline
Dimmock et al 2021 & FIN & 2017-09-07 &  &  & measured GIC (pipeline), \newline measured B, \newline calculated GIC, \newline calculated B & 1 site & http://space.fmi.fi/gic/ (FMI) & Pending \\ \hline
Wang et al 2021 & CHN & 2004-11-10 &  &  & calculated GIC & 34 sites & https://zenodo.org/record/4017356 (calculated) & Pending \\ \hline
Marsal et al 2021 & ESP & 2020 - 2021 & $\checkmark$ &  & measured B, \newline calculated GIC & 4 sites & https://doi.org/10.20350/digitalCSIC/14004 & Pending \\ \hline
Mukhtar et al 2020 & NZL & 2015-03-17, \newline 2003-11-20 &  &  & measured B, \newline calculated GIC & 30 sites & NA & No \\ \hline
Kelbert et al 2020 & USA & 2003-10-31 &  &  & measured B, \newline calculated GIC & 4 sites & NA & No \\ \hline
Divett et al 2020 & NZL & 2015-03-17 &  & $\checkmark$ & measured B, \newline calculated GIC & 23 sites & NA & No \\ \hline
Zhang et al 2020 & CHN & 2015-12-14, \newline 2016-03-06, \newline 2015-03-17 &  & $\checkmark$ & measured B, \newline calculated GIC & 5 sites & https://doi.org/10.5281/zenodo.3714269 & Pending \\ \hline
Heyns et al 2020 & RSA/USA & 2003-10-31, \newline 2015-06-23, \newline 1989-03-15 &  & $\checkmark$ & measured B, \newline calculated GIC & 2 sites & NA (TVA data shared via email) & Yes \\ \hline
Alves Ribeiro et al 2020 & PRT & 2015 &  &  & measured B, \newline calculated GIC & 20 sites & NA & No \\ \hline
Caraballo et al 2019 & MEX & 2000-07-15, \newline 2003-10-20, \newline 2015-03-17, \newline 2017-09-07 &  &  & measured B, \newline calculated GIC & 115 sites & https://doi.org/10.17605/OSF.IO/AV9HT "GIC\_assess\_v2:SW Journal Article 2020/datasets" (calculated) & Pending \\ \hline
Marsal et al 2019 & ESP & 2012-01-22 &  & $\checkmark$ & calculated GIC & 1 site & NA & No \\ \hline
Dimmock et al 2019 & FIN & 2017-09-07 &  &  & measured GIC (pipeline), \newline calculated GIC & 1 site & http://space.fmi.fi/gic/man\_ascii/man.php (FMI) & Pending \\ \hline
Marshall et al 2019 & AUS & 2013-03-17, \newline 2013-10-02, \newline 2015-06-23, \newline 2017-09-08 &  & $\checkmark$ & measured B, \newline calculated GIC & 8 sites & NA & No \\ \hline
Khanal et al 2019 & FIN & 113 events (list in SI of paper) &  &  & measured GIC (pipeline), \newline measured B & 1 site & http://space.fmi.fi/gic/?page=gasum\_final (FMI), \newline https://agupubs.onlinelibrary.wiley.com/action/downloadSupplement?doi=10.1029\%2F2018SW001879\&file=swe20813-sup-0002-2018SW001879-ds01.zip (datasets used in paper) & Pending \\ \hline
Rosenqvist et al 2019 & SWE & 2015-07-05 &  & $\checkmark$ & measured B, \newline calculated GIC & 1 site & NA & No \\ \hline
Espnosa et al 2019 & BRA & 2013-10-08, \newline 2015-03-17, \newline 2015-06-22 &  &  & measured GIC (2013 storm only), \newline measured B, \newline calculated GIC & 1 site & https://agupubs.onlinelibrary.wiley.com/action/downloadSupplement?doi=10.1029\%2F2018SW002094\&file=swe20818-sup-0002-2018SW002094-Table\_SI-S01.xlsx & Pending \\ \hline
Weigel et al 2019 & JPN & 2006 - 20007 &  & $\checkmark$ & measured B, \newline calculated GIC & 1 site & NA (data from Watari shared via email) & Pending \\ \hline
Tozzi et al 2018 & ITA & 1999-2015 &  &  & measured B, \newline calculated GIC index & 2 sites & NA & No \\ \hline
Nakamura et al 2018 & JPN & 2017-05-27 &  & $\checkmark$ & measured B, \newline calculated GIC & 2 sites & NA & No \\ \hline
Blake et al 2018 & IRL & 2015-08-26, \newline 2015-09-07, \newline 2015-10-07, \newline 2015-12-20, \newline 2016-03-06 &  & $\checkmark$ & measured B, \newline calculated GIC & 1 site & https://agupubs.onlinelibrary.wiley.com/action/downloadSupplement?doi=10.1029\%2F2018SW001926\&file=swe20785-sup-0002-2018SW001926-Data\_Set\_SI-S01.txt (2015-08), \newline https://agupubs.onlinelibrary.wiley.com/action/downloadSupplement?doi=10.1029\%2F2018SW001926\&file=swe20785-sup-0003-2018SW001926-Data\_Set\_SI-S02.txt (2015-09), \newline https://agupubs.onlinelibrary.wiley.com/action/downloadSupplement?doi=10.1029\%2F2018SW001926\&file=swe20785-sup-0004-2018SW001926-Data\_Set\_SI-S03.txt (2015-10), \newline https://agupubs.onlinelibrary.wiley.com/action/downloadSupplement?doi=10.1029\%2F2018SW001926\&file=swe20785-sup-0005-2018SW001926-Data\_Set\_SI-S04.txt (2015-12), \newline https://agupubs.onlinelibrary.wiley.com/action/downloadSupplement?doi=10.1029\%2F2018SW001926\&file=swe20785-sup-0006-2018SW001926-Data\_Set\_SI-S05.txt (2016-03) & Pending \\ \hline
Bailey et al 2018 & AUT & 2017-04-23 &  & $\checkmark$ & measured B, \newline calculated GIC & 3 sites & NA (paper has map of AUT system) & No \\ \hline
Clilverd et al 2018 & NZL & 2017-09-07 &  & $\checkmark$ & measured B, \newline calculated GIC & 1 site & NA & No \\ \hline
Divett et al 2018 & NZL & generalized event &  &  & measured B, \newline calculated GIC & 29 sites & NA & No \\ \hline
Mac Manus et al 2017 & NZL & 2001-11-06, \newline 2013-10-02 &  & $\checkmark$ & measured B & 17 sites & NA & No \\ \hline
Torta et al 2017 & ESP & 2003-10-29, \newline 2015-03-17 &  & $\checkmark$ & measured B, \newline calculated GIC & 1 site & NA & No \\ \hline
Barbosa et al 2017 & BRA/RSA/GBR/FIN & 90 storms &  &  & measured B, \newline calculated GIC & 4 sites & NA & No \\ \hline
Divett et al 2017 & NZL & generalized event &  &  & calculated GIC & 29 sites & NA & No \\ \hline
Butala et al 2017 & USA & 2013-10-02, \newline 2013-06-01, \newline 2013-10-09 &  & $\checkmark$ & measured B, \newline calculated GIC & 4 sites & https://agupubs.onlinelibrary.wiley.com/action/downloadSupplement?doi=10.1002\%2F2017SW001602\&file=swe20522-sup-0002-supplementary.zip (daily plots only) & Pending \\ \hline
Gannon et al 2017 & USA & 1989-03-13, \newline 1991-03-24, \newline 2003-10-28, \newline 2016-04-12, \newline 2016-05-07 &  &  & measured B & calculated GIC & NA & No \\ \hline
Rodger et al 2017 & NZL & 2001 - 2015 &  & $\checkmark$ & measured B, \newline calculated GIC & 2 sites & NA & No \\ \hline
Ingham et al 2017 & NZL & 2015-03-17, \newline 2013-06-01, \newline 2013-06-29, \newline 2013-10-02, \newline 2014-09-12 &  & $\checkmark$ & measured B, \newline calculated GIC & 5 sites & NA & No \\ \hline
Blake et al 2016 & IRL & 2016-03-06, \newline 2015-12-20, \newline 2015-03-17, \newline 2003-10-29, \newline 1989-03-13 &  & $\checkmark$ & measured B, \newline calculated GIC & 1 site & NA & No \\ \hline
Liu et al 2016 & CHN & 2015-03-17, \newline 2015-06-22 &  &  & measured GIC (railway), \newline measured B & 1 site & NA & NA \\ \hline
Kelly et al 2016 & GBR/FRA & 1989-03-13, \newline 2003-10-29, \newline 1991-06-09 &  &  & measured B, \newline calculated GIC & 4 sites & NA & No \\ \hline
Boteler et al 2016 & generalized & generalized event &  &  & calculated GIC & 1 site & NA & No \\ \hline
Zhang et al 2016 & CHN & 2012-07-23 &  &  & calculated B, \newline calculated GIC indices & 1 site & NA & No \\ \hline
Liu et al 2016 & CHN & 2012-07-14 &  & $\checkmark$ & measured B & 3 sites & NA & No \\ \hline
Matandirotya et al 2016 & RSA & 2013-07-06, \newline 2015-03-17 & $\checkmark$ &  & measured B, \newline calculated GIC & 2 sites & NA & No \\ \hline
Zhang et al 2015 & CHN & 2015-03-17 &  & $\checkmark$ & measured B, \newline calculated B, \newline calculated GIC & 2 sites & NA & No \\ \hline
\end{longtable}
%\caption{List of GIC papers and their data contents.}
%\label{table:gic_papers}
%\end{table}


%%%%%%%%%%%%%%%%%%%%%%%%%%%%%%%%%%%%%%%%%%%%%%%
% REFERENCES and BIBLIOGRAPHY
%
% \bibliography{<name of your .bib file>} don't specify the file extension
% don't specify bibliographystyle
%
%%%%%%%%%%%%%%%%%%%%%%%%%%%%%%%%%%%%%%%%%%%%%%%

%\bibliography{ enter your bibtex bibliography filename here }



%Reference citation instructions and examples:
%
% Please use ONLY \cite and \citeA for reference citations.
% \cite for parenthetical references
% ...as shown in recent studies (Simpson et al., 2019)
% \citeA for in-text citations
% ...Simpson et al. (2019) have shown...
%
%
%...as shown by \citeA{jskilby}.
%...as shown by \citeA{lewin76}, \citeA{carson86}, \citeA{bartoldy02}, and \citeA{rinaldi03}.
%...has been shown \cite{jskilbye}.
%...has been shown \cite{lewin76,carson86,bartoldy02,rinaldi03}.
%... \cite <i.e.>[]{lewin76,carson86,bartoldy02,rinaldi03}.
%...has been shown by \cite <e.g.,>[and others]{lewin76}.
%
% apacite uses < > for prenotes and [ ] for postnotes
% DO NOT use other cite commands (e.g., \citet, \citep, \citeyear, \nocite, \citealp, etc.).
%

\bibliography{main}

\end{document}



More Information and Advice:

%%%%%%%%%%%%%%%%%%%%%%%%%%%%%%%%%%%%%%%%%%%%%%%
%
%  SECTION HEADS
%
%%%%%%%%%%%%%%%%%%%%%%%%%%%%%%%%%%%%%%%%%%%%%%%

% Capitalize the first letter of each word (except for
% prepositions, conjunctions, and articles that are
% three or fewer letters).

% AGU follows standard outline style; therefore, there cannot be a section 1 without
% a section 2, or a section 2.3.1 without a section 2.3.2.
% Please make sure your section numbers are balanced.
% ---------------
% Level 1 head
%
% Use the \section{} command to identify level 1 heads;
% type the appropriate head wording between the curly
% brackets, as shown below.
%
%An example:
%\section{Level 1 Head: Introduction}
%
% ---------------
% Level 2 head
%
% Use the \subsection{} command to identify level 2 heads.
%An example:
%\subsection{Level 2 Head}
%
% ---------------
% Level 3 head
%
% Use the \subsubsection{} command to identify level 3 heads
%An example:
%\subsubsection{Level 3 Head}
%
%---------------
% Level 4 head
%
% Use the \subsubsubsection{} command to identify level 3 heads
% An example:
%\subsubsubsection{Level 4 Head} An example.
%
%%%%%%%%%%%%%%%%%%%%%%%%%%%%%%%%%%%%%%%%%%%%%%%
%
%  IN-TEXT LISTS
%
%%%%%%%%%%%%%%%%%%%%%%%%%%%%%%%%%%%%%%%%%%%%%%%
%
% Do not use bulleted lists; enumerated lists are okay.
% \begin{enumerate}
% \item
% \item
% \item
% \end{enumerate}
%
%%%%%%%%%%%%%%%%%%%%%%%%%%%%%%%%%%%%%%%%%%%%%%%
%
%  EQUATIONS
%
%%%%%%%%%%%%%%%%%%%%%%%%%%%%%%%%%%%%%%%%%%%%%%%

% Single-line equations are centered.
% Equation arrays will appear left-aligned.

Math coded inside display math mode \[ ...\]
 will not be numbered, e.g.,:
 \[ x^2=y^2 + z^2\]

 Math coded inside \begin{equation} and \end{equation} will
 be automatically numbered, e.g.,:
 \begin{equation}
 x^2=y^2 + z^2
 \end{equation}


% To create multiline equations, use the
% \begin{eqnarray} and \end{eqnarray} environment
% as demonstrated below.
\begin{eqnarray}
  x_{1} & = & (x - x_{0}) \cos \Theta \nonumber \\
        && + (y - y_{0}) \sin \Theta  \nonumber \\
  y_{1} & = & -(x - x_{0}) \sin \Theta \nonumber \\
        && + (y - y_{0}) \cos \Theta.
\end{eqnarray}

%If you don't want an equation number, use the star form:
%\begin{eqnarray*}...\end{eqnarray*}

% Break each line at a sign of operation
% (+, -, etc.) if possible, with the sign of operation
% on the new line.

% Indent second and subsequent lines to align with
% the first character following the equal sign on the
% first line.

% Use an \hspace{} command to insert horizontal space
% into your equation if necessary. Place an appropriate
% unit of measure between the curly braces, e.g.
% \hspace{1in}; you may have to experiment to achieve
% the correct amount of space.


%%%%%%%%%%%%%%%%%%%%%%%%%%%%%%%%%%%%%%%%%%%%%%%
%
%  EQUATION NUMBERING: COUNTER
%
%%%%%%%%%%%%%%%%%%%%%%%%%%%%%%%%%%%%%%%%%%%%%%%

% You may change equation numbering by resetting
% the equation counter or by explicitly numbering
% an equation.

% To explicitly number an equation, type \eqnum{}
% (with the desired number between the brackets)
% after the \begin{equation} or \begin{eqnarray}
% command.  The \eqnum{} command will affect only
% the equation it appears with; LaTeX will number
% any equations appearing later in the manuscript
% according to the equation counter.
%

% If you have a multiline equation that needs only
% one equation number, use a \nonumber command in
% front of the double backslashes (\\) as shown in
% the multiline equation above.

% If you are using line numbers, remember to surround
% equations with \begin{linenomath*}...\end{linenomath*}

%  To add line numbers to lines in equations:
%  \begin{linenomath*}
%  \begin{equation}
%  \end{equation}
%  \end{linenomath*}



